% !TEX root = chapter-standalone.tex

\section{Powerset}
\SYNDEF{powerset}
\label{sec:power-set}

\begin{ctdefinition}[Power set]
    \label{def:power-set}
    Given a set~\setA, we can form a new set whose elements are precisely all the subsets of~\setA.
    This new set is called the \emph{powerset} of~\setA; we denote it by~$\powerset \setA$.
\end{ctdefinition}

For example, if~$\setA = \makeset{ \sapple, \scarrot, \sgrapes }$, then its \SY{powerset} is
\begin{equation}
    \powerset \setA = \makeset{ \Emptyset, \makeset{ \sapple }, \makeset{ \scarrot }, \makeset{ \sgrapes}, \makeset{ \sapple, \scarrot }, \makeset{ \scarrot, \sgrapes }, \makeset{ \sapple, \sgrapes}, \makeset{ \sapple, \scarrot, \sgrapes } }.
\end{equation}

\begin{exercise}
    Can you count how many elements the \SY{powerset}~$\powerset \setA$ has in the following cases?
    \begin{enumerate}
        \item $\setA = \makeset{ \sapple } $.
        \item $\setA = \makeset{ \sapple, \scarrot } $.
        \item $\setA = \makeset{ \sapple, \scarrot, \sgrapes } $.
        \item $\setA = \Emptyset $.
    \end{enumerate}
    Can you guess a general formula for the size of the \SY{powerset} of a finite set?
\end{exercise}

\begin{solution}
    We have:
    \begin{enumerate}
        \item 2.
        \item 4.
        \item 8.
        \item 1.
    \end{enumerate}
    In general, the size of~$\powerset \setA$ is~$2$ to the power of the size of~\setA.
\end{solution}

Now suppose we fix a set~\setA for a moment.
Given~$\subA \setin\powerset \setA$, the \emph{complement} of~$\subA$ with respect to~\setA is
\begin{equation}
    \setA \setcomplement \subA = \makeset{ \ela \setin\setA \mid \ela \notsetin \subA },
\end{equation}
which is again an element of~$\powerset \setA$.
In situations where it is evident which ambient set~\setA we are working with, the notation~$\compl{\subA}$ is sometimes used instead of~$\setA \setcomplement \subA$.

We also note that the operations of union and intersection, when restricted to~$\powerset \setA$, again produce elements of~$\powerset \setA$.
That is, if~$\subA, \subB \setin\powerset \setA$, then~$\subA \setunion \subB \setin\powerset \setA$, and similarly~$\subA \setintersection \subB \setin\powerset \setA$.

The operations~$\setunion$,~$\setintersection$, and~$\compl{( - )}$ on~$\powerset \setA$ obey various rules which are useful to be familiar with.
For example,
\begin{equation}
    \compl{(\subA \setintersection \subB)} = (\compl{\subA}) \setunion (\compl{\subB}).
\end{equation}
Can you state more such rules?
A useful visual aid for such calculations are so-called Venn diagrams (\cref{fig:venn-union}, \cref{fig:venn-intersection}, \cref{fig:venn-complement}).

\begin{marginfigure}
    \centering
    \includesag{venn-union}
    \caption{Venn diagram for union operation.}
    \label{fig:venn-union}
\end{marginfigure}

\begin{marginfigure}
    \centering
    \includesag{venn-intersection}
    \caption{Venn diagram for intersection operation.}
    \label{fig:venn-intersection}
\end{marginfigure}

\begin{marginfigure}
    \centering
    \includesag{venn-complement}
    \caption{Venn diagram for complement operation.}
    \label{fig:venn-complement}
\end{marginfigure}

\begin{gradedexercise}[\exname{DistributingSubsets}]
    \label{ex:distributing-subsets}
    Let~\setA be a set, and let~$\subA, \subB, \subC \setsubseteq \setA$ be subsets.
    Prove that
    \begin{equation}
        \subA \setintersection (\subB \setunion \subC) = (\subA \setintersection \subB) \setunion (\subA \setintersection \subC).
    \end{equation}
\end{gradedexercise}

\solutionof{DistributingSubsets}






